\section{Where Classicality Lives: the R3 Hinge}
\label{sec:hinge}

Classicality lives at multiple levels in this work: one
meta-theoretic, in the proof assistant itself, and one object-level,
in the frame's own logic.

\paragraph{Lean's meta-theory.}
Lean's foundations include
\texttt{Classical.choice}, and \textsf{mathlib} is built on it
throughout.  For example, the order relation on 
$\mathbb{R}$ and $\ENNReal$ has no
computable \texttt{Decidable} instance in general, so any library
lemma that case-splits on $x \le y$ for reals resolves that split via
\texttt{Classical.propDecidable}. A
formalization that never invokes excluded middle at the level of
frame elements $a \in \Omega$ can therefore still be classical 
through the arithmetic it borrows from the library.  We track
this separately from object-level classicality throughout the paper.
Theorem~\ref{thm:decomposition} of Section~\ref{sec:representation}
rests only on this ambient baseline and contains no explicit case
split. The Scott-continuity recovery theorem
(Theorem~\ref{thm:recovery})
invokes Lean's \texttt{classical} tactic directly, and the
\texttt{CoxModel} witness of Section~\ref{sec:cox} goes the other
way, replacing the default with an explicit computable
\texttt{Decidable} instance to keep that one instance genuinely
non-classical.

\paragraph{Object-level classicality: the R3 axiom.}
Van Horn's repaired Cox theorem~\cite{vanhorn2003} axiomatizes
plausibility over classical propositions.  Its negation axiom, called
\emph{R3}, posits that the plausibility of $\neg A$ is a fixed
non-increasing function $S$ of the plausibility of $A$.  Consistency
forces the involution $S \circ S = \mathrm{id}$, and therefore the
complement rule $p \mapsto 1 - p$.  We show that the 
involutive negation \emph{is}
double-negation elimination, so R3 is where classical logic enters
every known Cox-style derivation.  To do this, we first 
isolate the exact content of
R3's conclusion as a predicate on valuations, then prove equivalence
to excluded middle.

\begin{lstlisting}
def HasClassicalNegation (v : Valuation Ω) : Prop :=
  ∀ a, v a + v aᶜ = 1
\end{lstlisting}

\begin{theorem}[the R3 hinge]
\label{thm:hinge}
For every frame $\Omega$, every valuation on $\Omega$ has classical
negation \emph{iff} $a \sqcup a^{\mathsf{c}} = \top$ for every
$a \in \Omega$.
\end{theorem}

The forward direction (\texttt{hasClassicalNegation\_of\_em}) is a
calculation.  The content is the converse
(\texttt{em\_of\_forall\_\allowbreak hasClassicalNegation}): wherever excluded
middle fails, a valuation with genuine slack must be
\emph{manufactured}.  Suppose $a \sqcup a^{\mathsf{c}} \neq \top$.
The prime ideal separation theorem
(\textsf{mathlib}'s \texttt{DistribLattice.\allowbreak prime\_ideal\_of\_%
disjoint\_filter\_ideal}) yields a prime ideal $J$ containing
$a \sqcup a^{\mathsf{c}}$ but not $\top$.  The indicator of $J$'s
complement,
\[
  v(x) = \begin{cases} 0 & x \in J \\ 1 & x \notin J, \end{cases}
\]
is a valuation.  Modularity in the only nontrivial case ($x, y
\notin J$) is \emph{exactly} primeness, $x \sqcap y \in J
\Rightarrow x \in J \vee y \in J$.  And this valuation assigns
$v(a) + v(a^{\mathsf{c}}) = 0 \neq 1$.  So the complement rule is not
one convention among many.  Quantified over all valuations, it is a
statement \emph{about the algebra} that is equivalent to excluded
middle.

\paragraph{Sharp valuations are the points of the locale.}
The construction above is not ad hoc.  It is one half of an exact
correspondence.  Call $v$ \emph{sharp} if it is $\{0,1\}$-valued, a
state of full certainty.  Sharpness characterizes exactly the
prime-ideal indicators:

\begin{theorem}[\texttt{sharp\_iff\_point}]
For every valuation $v$ on a frame $\Omega$, $v$ is sharp iff $v$ is
the complement-indicator of a prime ideal.
\end{theorem}

\noindent
This theorem is witnessed by two mutually inverse maps: $v \mapsto Z(v) =
\{x \in \Omega : v(x) = 0\}$, which computes the zero-set of a sharp valuation, 
and the map taking
$J$ to its complement-indicator (the valuation of
Theorem~\ref{thm:hinge}'s construction).  $Z(v)$ is prime exactly
because $v$ is sharp, and the complement-indicator of $Z(v)$
recovers $v$ itself.  In locale-theoretic terms, sharp valuations
recover the \emph{points} of the locale.  A prime ideal is
\emph{completely prime} when it is also closed under arbitrary
(not just finite) suprema of its members. 
The finitely prime points 
of \texttt{sharp\_iff\_point} sharpen to
the completely prime points exactly where Scott continuity is added,
recovering points of the locale in the sense used
by~\cite{vickers2008,coquand-spitters2009}:

\begin{theorem}[\texttt{sharp\_scott\_iff\_\allowbreak completelyPrimePoint}]
For every valuation $v$ on a frame $\Omega$, $v$ is sharp and
Scott-continuous iff $v$ is the complement-indicator of a completely
prime ideal.
\end{theorem}

\noindent  The hinge and the two
correspondences together answer the question ``where
does classicality live'': whenever
excluded middle fails at some $a$, the failure is witnessed by a
point of the locale.
